\chapter[Chapter \thechapter]{}
\lettrine{T}{he} scramble for places subsided; the jury returned; the prisoner reappeared in the dock suddenly, like a jack-in-the-box; the judge resumed his seat. Some petals had spilt from the roses. The old voice took up its tale where it had left off.

<Members of the jury—there is no need, I think, for me to recall the course of Philip Boyes' illness in great detail. The nurse was called in on June 21st, and during that day the doctors visited the patient three times. His condition grew steadily worse. There was persistent vomiting and diarrhoea, and he could not keep any food or medicine down at all. On the day after, the 22nd, he was worse still—in great pain, the pulse growing weaker, and the skin about the mouth getting dry and peeling off. The doctors gave him every attention, but could do nothing for him. His father was summoned, and when he arrived he found his son conscious, but unable to lift himself. He was able to speak, however, and in the presence of his father and Nurse Williams he made the remark, <I'm going out, Dad, and I'm glad to be through with it. Harriet'll be rid of me now—I didn't know she hated me quite so much.> Now that was a very remarkable speech, and we have heard two very different interpretations put upon it. It is for you to say whether, in your opinion, he meant: <She has succeeded in getting rid of me; I didn't know she hated me enough to poison me,> or whether he meant, <When I realised she hated me so much, I decided I did not want to live any longer>—or whether, perhaps, he meant neither of these things. When people are very ill, they sometimes get fantastic ideas, and sometimes they wander in their minds; perhaps you may feel that it is not profitable to take too much for granted. Still, those words are part of the evidence, and you are entitled to take them into account.

During the night he became gradually weaker and lost consciousness, and at 3 o'clock in the morning he died, without ever regaining it. That was on the 23rd. of June.

Now, up to this time, no suspicion of any kind had been aroused. Both Dr~Grainger and Dr~Weare formed the opinion that the cause of death was acute gastritis, and we need not blame them for coming to this conclusion, because it was quite consistent both with the symptoms of the illness and with the past history of the patient. A death-certificate was given in the usual way, and the funeral took place on the 28th.

Well, then something happened which frequently does happen in cases of this kind, and that is that somebody begins to talk. It was Nurse Williams who talked in this particular case, and while you will probably think that this was a very wrong and a very indiscreet thing for a nurse to do, yet, as it turns out, it was a good thing that she did. Of course, she ought to have told Dr~Weare or Dr~Grainger of her suspicions at the time, but she did not do this, and we may at least feel glad to know that, in the doctors' opinions, even if she had done so, and if they had discovered that the illness was caused by arsenic, they would not have been able to do anything more to save the life of this unfortunate man. At any rate, what happened was that Nurse Williams was sent, during the last week of June, to nurse another patient of Dr~Weare's, who happened to belong to the same literary set in Bloomsbury as Philip Boyes and Harriet Vane, and while she was there, she spoke about Philip Boyes, and said that, in her opinion, the illness looked very much like poisoning, and she even mentioned the word arsenic. Well, you know how a thing like that gets about. One person tells another and it is discussed at tea-parties, or what are known, I believe, as cocktail parties, and very soon a story gets spread about, and people mention names and take sides. Miss~Marriott and Miss~Price were told about it, and it also got to the ears of Mr~Vaughan. Now Mr~Vaughan had been greatly distressed and surprised by Philip Boyes' death, especially as he had been with him in Wales, and knew how much he had improved in health while on his holiday, and he also felt very strongly that Harriet Vane had behaved badly about the love-affair. Mr~Vaughan felt that some action ought to be taken about the matter, and went to Mr~Urquhart and put the story before him. Now Mr~Urquhart is a solicitor, and is therefore inclined to take a cautious view of rumours and suspicions, and he warned Mr~Vaughan that it was not wise to go about making accusations against people, for fear of an action for libel. At the same time, he naturally felt uneasy that such a thing should be said about a relation who had died in his house. He took the course—the very sensible course—of consulting Dr~Weare and suggesting that, if he was quite certain that the illness was due to gastritis and nothing else, he should take steps to rebuke Nurse Williams and put an end to the talk. Dr~Weare was naturally very much surprised and upset to hear what was being said, but, since the suggestion had been made, he could not deny that—taking the symptoms only into account—there was just the bare possibility of something of the sort, because, as you have already heard in the medical evidence, the symptoms of arsenical poisoning and of acute gastritis are really indistinguishable.

When this was communicated to Mr~Vaughan, he was confirmed in his suspicions, and wrote to the elder Mr~Boyes suggesting an enquiry. Mr~Boyes was naturally very much shocked, and said at once that the matter should be taken up. He had known of the liaison with Harriet Vane, and had noticed that she did not come to enquire after Philip Boyes, nor attend the funeral, and this had struck him as heartless behaviour. In the end, the police were communicated with and an exhumation order obtained.

You have heard the result of the analysis made by Sir~James Lubbock and Mr~Stephen Fordyce. There was a great deal of discussion about methods of analysis and the way that arsenic behaves in the body and so on, but, I think we need not trouble too much about those fine details. The chief points in the evidence seemed to me to be these, which you may note down if you care to do so.

The analysts took certain organs of the body—the stomach, intestines, kidneys, liver and so on, and analysed portions of these and found that they all contained arsenic. They were able to weigh the quantity of arsenic found in these various portions, and they calculated from that the quantity of arsenic present in the whole body. Then they had to allow so much for the amount of arsenic eliminated from the body by the vomiting and diarrhoea and also through the kidneys, because the kidneys play a very large part in the elimination of this particular poison. After making allowance for all these things, they formed the opinion that a large and fatal dose of arsenic—four or five grains, perhaps, had been taken about three days before the death.

I do not know whether you quite followed all the technical arguments about this. I will try to tell you the chief points as I understood them. The nature of arsenic is to pass through the body very quickly, especially if it is taken with food or immediately following a meal, because the arsenic irritates the lining of the internal organs and speeds up the process of elimination. The action would be quicker if the arsenic were taken in liquid than if it were taken in the form of a powder. Where arsenic was taken with, or immediately on top of a meal, nearly the whole of it would be evacuated within twenty-four hours after the onset of the illness. So you see that, although the actual quantities found in the body may seem to you and me very small indeed, the mere fact that they were found there at all, after three days of persistent vomiting and diarrhoea and so on, points to a large dose having been taken at some time.

Now there was a great deal of discussion about the time at which the symptoms first set in. It is suggested by the defence that Philip Boyes may have taken the arsenic himself at some time between leaving Harriet Vane's flat and hailing the taxi in Guilford Street; and they bring forward books which show that in many cases the onset of symptoms takes place in a very short time after taking the arsenic—a quarter of an hour, I think, was the shortest time mentioned where the arsenic was taken in liquid form. Now the prisoner's statement—and we have no other—is that Philip Boyes left her at 10 o'clock, and at ten minutes past he was in Guilford Street. He was then looking ill. It would not take many minutes to drive to Woburn Square at that hour of night, and by the time he got there, he was already in acute pain and hardly able to stand. Now Guilford Street is a very short way from Doughty Street—perhaps three minutes' walk—and you must ask yourselves, if the prisoner's statement is correct, what he did with those ten minutes. Did he occupy himself in going to some quiet spot and taking a dose of arsenic, which he must in that case have brought with him in anticipation of an unfavourable interview with the prisoner? And I may remind you here, that the defence have brought no evidence to show that Philip Boyes ever bought any arsenic, or had access to any arsenic. That is not to say he could not have obtained it—the purchases made by Harriet Vane show that the law about the sale of poisons is not always as effective as one would like it to be—but the fact remains that the defence have not been able to show that the deceased ever had arsenic in his possession. And while we are on this subject, I will mention that, curiously enough, the analysts could find no traces of the charcoal, or indigo, with which commercial arsenic is supposed to be mixed. Whether it was bought by the prisoner or by the deceased himself, you would expect to find traces of the colouring matter. But you may think it likely that all such traces would be removed from the body by the vomiting and purging which took place.

As regards the suggestion of suicide, you will have to ask yourselves about those ten minutes—whether Boyes was taking a dose of arsenic, or whether, as is also possible, he felt unwell and sat down somewhere to recover himself, or whether, perhaps, he was merely roaming about in the vague way we sometimes do when we are feeling upset and unhappy. Or you may think that the prisoner was mistaken, or not speaking the truth, about the time he left the flat.

You have also the prisoner's statement that Boyes mentioned, before he left her, that he was feeling unwell. If you think this had anything to do with the arsenic, it of course disposes of the suggestion that he took poison after leaving the flat.

Then, when one looks into it, one finds that this question about the onset of symptoms is left very vague. Various doctors came here and told you about their own experiences and the cases quoted by medical authorities in books, and you will have noticed that there is no certainty at all about the time when the symptoms may be expected to appear. Sometimes it is a quarter of an hour or half an hour, sometimes two hours, sometimes as much as five or six, and, I believe, in one case as much as seven hours after taking the poison.>

Here the Attorney-General rose respectfully and said: <In that case, me lud, I think I am right in saying that the poison was taken on an empty stomach.>

<Thank you, I am much obliged to you for the reminder. That was a case in which the poison was taken on an empty stomach. I only mention these cases to show that we are dealing with a very uncertain phenomenon, and that is why I was particular to remind you of all the occasions on which Philip Boyes took food during the day—the 20th of June, since there is always the possibility that you may have to take them into consideration.>

<A beast, but a just beast,> murmured Lord~Peter Wimsey.

<I have purposely left out of consideration until now another point which arose out of the analysis, and that is the presence of arsenic in the hair. The deceased had curly hair, which he wore rather long; the front portion, when straightened out, measured about six or seven inches in places. Now, in this hair, arsenic was found, at the end closest to the head. It did not extend to the tips of the longest hair, but it was found near the roots, and Sir~James Lubbock says that the quantity was greater than could be accounted for in any natural way. Occasionally, quite normal people are found to have minute traces of arsenic in the hair and skin and so on, but not to the amount found here. That is Sir~James' opinion.

Now you have been told—and the medical witnesses all agree in this—that if a person takes arsenic, a certain proportion of it will be deposited in the skin, nails and hair. It will be deposited in the root of the hair, and as the hair grows, the arsenic will be carried along with the growth of the hair, so that you get a rough idea, from seeing the position of the arsenic in the hair, how long the administration of arsenic has been going on. There was a good deal of discussion about this, but I think there was a fairly general agreement that, if you took a dose of arsenic, you might expect to find traces of it in the hair, close to the scalp, after about ten weeks. Hair grows at the rate of about six inches in a year, and the arsenic will grow out with it till it reaches the far end and is cut off. I am sure that the ladies on the jury will understand this very well, because I believe that the same thing occurs in the case of what is termed a <permanent wave.> The wave is made in a certain portion of the hair, and after a time it grows out, and the hair near the scalp comes up straight and has to be waved again. You can tell by the position of the wave, how long ago the waving was done. In the same way, if a finger-nail is bruised, the discolouration will gradually grow up the nail until it reaches the point where you can cut it off with the scissors.

Now it has been said that the presence of arsenic in and about the roots of Philip Boyes' hair indicates that he must have taken arsenic three months at least before his death. You will consider what importance is to be attached to this in view of the prisoner's purchases of arsenic in April and May, and of the deceased's attacks of sickness in March, April and May. The quarrel with the prisoner took place in February; he was ill in March and he died in June. There are five months between the quarrel and the death, and four months between the first illness and the death, and you may think that there is some significance in these dates.

We now come to the enquiries made by the police. When suspicion was aroused, detectives investigated Harriet Vane's movements and subsequently went to her flat to take a statement from her. When they told her that Boyes was found to have died of arsenic poisoning, she appeared very much surprised, and said, <Arsenic? What an extraordinary thing!> And then, she laughed, and said, <Why, I am writing a book all about arsenic poisoning.> They asked her about the purchases of arsenic and other poisons which she had made and she admitted them quite readily and at once gave the same explanation that she gave here in court. They asked what she had done with the poisons, and she replied that she had burnt them because they were dangerous things to have about. The flat was searched, but no poisons of any kind were found, except such things as aspirin and a few ordinary medicines of that kind. She absolutely denied having administered arsenic or any kind of poison to Philip Boyes. She was asked whether the arsenic could possibly have got into the coffee by accident, and replied that that was quite impossible, as she had destroyed all the poisons before the end of May.>

Here Sir~Impey Biggs interposed and begged with submission to suggest that his lordship should remind the jury of the evidence given by Mr~Challoner.

<Certainly, Sir~Impey, I am obliged to you. You remember that Mr~Challoner is Harriet Vane's literary agent. He came here to tell us that he had discussed with her as long ago as last December the subject of her forthcoming book, and she then told him that it was to be about poisons, and very probably about arsenic. So you may think it is a point in the prisoner's favour that this intention of studying the purchase and administration of arsenic was already in her mind some time before the quarrel with Philip Boyes took place. She evidently gave considerable thought to the subject, for there were a number of books on her shelves dealing with forensic medicine and toxicology, and also the reports of several famous poison trials, including the Madeleine Smith case, the Seddon case and the Armstrong case—all of which were cases of arsenical poisoning.

Well, I think that is the case as it is presented to you. This woman is charged with having murdered her former lover by arsenic. He undoubtedly did take arsenic, and if you are satisfied that she gave it to him with intent to injure or kill him, and that he died of it, then it is your duty to find her guilty of murder.

Sir~Impey Biggs, in his able and eloquent speech, has put it to you that she had very little motive for such a murder, but I am bound to tell you that murders are very often committed for what seem to be most inadequate motives—if, indeed, any motive can be called adequate for such a crime. Especially where the parties are husband and wife, or have lived together as husband and wife, there are likely to be passionate feelings which may tend to crimes of violence in persons with inadequate moral standards and of unbalanced mind.

The prisoner had the means—the arsenic—she had the expert knowledge, and she had the opportunity to administer it. The defence say that this is not enough. They say the Crown must go further and prove that the poison could not have been taken in any other way—by accident, or with suicidal intent. That is for you to judge. If you feel that there is any reasonable doubt that the prisoner gave this poison to Philip Boyes deliberately, you must bring her in Not Guilty of murder. You are not bound to decide how it was given, if it was not given by her. Consider the circumstances of the case as a whole, and say what conclusion you have come to.>